\section{TWAP Design}

TWAP is implemented as an absolute schedule, not a sleep loop. At monotonic elapsed time \(t\):

\begin{equation}
  scheduled\_cumulative\_quantity(t)
  = total\_trade\_quantity \times \frac{elapsed\_time}{total\_duration}.
\end{equation}

\begin{equation}
  quantity\_deficit(t)
  = scheduled\_cumulative\_quantity(t) -
    confirmed\_cumulative\_filled\_quantity(t).
\end{equation}

Before submitting, the engine subtracts reserved exposure:

\begin{equation}
\begin{split}
safe\_child\_quantity = quantity\_deficit
 - (&live\_open + pending\_submit \\
&+ pending\_cancel + unknown\_order).
\end{split}
\end{equation}

This design solves the problems that make the invalid \code{slice\_qty / sleep} implementation unsafe. Absolute schedule time avoids accumulated sleep drift. Unfilled earlier quantity naturally carries forward into the next deficit. Open or unknown earlier slices cannot be submitted twice because reserved exposure is subtracted. The final slice can absorb legal step-size rounding remainder, but the invariant still prevents exceeding the normalized target.

The scheduling window ends at \code{target\_duration\_seconds}. After that point the engine does not create new passive scheduled TWAP children. Deadline handling is a separate bounded cleanup phase. Under \code{CANCEL\_REMAINDER}, cleanup cancels and reconciles remaining exposure. Under \code{AGGRESSIVE\_WITHIN\_RANGE}, cleanup may first cancel passive exposure and then submit a marketable limit child, but only within the configured price bound and only if the exposure invariant permits it. The summary therefore reports \code{decision\_deadline\_elapsed\_seconds}, \code{terminal\_cleanup\_duration\_seconds}, \code{total\_lifecycle\_duration\_seconds}, and post-deadline order counts.

For example, suppose a \(0.010\) BTC buy is scheduled over 100 seconds. At 40 seconds the planned cumulative quantity is \(0.004\) BTC. If confirmed fills are only \(0.001\) BTC and an earlier child still reserves \(0.001\) BTC, the schedule deficit is \(0.003\) BTC but the safe new child quantity is \(0.002\) BTC. The missing earlier quantity is carried forward, while the still-open quantity is not submitted twice. This is the important difference between a schedule deficit and a child order size.

The deterministic TWAP artifact at \code{reports/evidence/simulation/twap/exec\_61fadac604f4440a} shows the same schedule behavior. Over a 100 second target, it submitted five 20-second slices of \(0.002\) BTC each, reconciled every slice fill, and completed exactly \(0.010\) BTC without exceeding the scheduled deficit.

The accepted Testnet TWAP run at \code{reports/evidence/testnet/twap/exec\_85bef3985ea3431a} also records the planned cumulative quantity, submitted quantity, open quantity, filled quantity, and unfilled quantity for each slice. Its first two slices were below the legal order shape and therefore submitted nothing; later slices carried the deficit forward, submitted legal child orders, and completed the required \(0.004\) BTC with \code{overfill\_quantity=0}.
